\textbf{Abstract.} The previous 64-residue construction guarantees that the direct gain \(g_0\) recovers at least \(10.004635224\%\) of a selected witness's odd-position allowance. We optimize the rates and the retained residues, with the numeric labeling allowed to be adversarial. For stationary 64-residue profiles of mean one, with realized allowed edges and score compatibility certified by the literal marked-tile window test, the sharp robust value is
\[
                \boxed{10.11902517624\%\ \text{(rounded)}}
\]
The previous rule already chooses a worst-case-optimal support when its rates are retuned. Allowing it to change supports helps on many individual labelings. It guarantees \(11\%\) on 15,318 of 16,384 normalized sparse configurations, uniformly in the dense completions. An explicit adverse order prevents any support policy in this class from improving the all-label value.

The upper proof reduces a large finite optimization to two scalar constraints. At the adverse order, retain every legal structural support and relax each rate up to \(2.39213\). An exhaustive calculation leaves only two supports whose relaxed value can exceed \(10.118\%\), once the necessary root mark is imposed. Both require the same sparse singleton and the same count adjustment. One pays an additional eight denominator units per common saturated rate. The other has the optimized coefficients of the earlier construction. The lower certificate and the two upper calculations consequently meet in the same analytic expression. Exact rational brackets put that value between \(10.1190251762367\%\) and \(10.1190251762403\%\). At one preassigned source exponent the optimum can be slightly higher; we give its parameter-dependent formula.

The arithmetic interpretation concerns possible retained sets passing the original score on a fixed conductor-one host, with the same set used for all its primes and lifts. The proof's original chosen set remains existential. A Lean replay checks two finite algebraic declarations. The support exhaustions have independently implemented executable checks; the transport to the original families is a mathematical argument with pinned source references. The investigation derives no new conclusion about the zeros of the zeta function.
